\section{Магнитное кольцо: второй порядок по фазе}
\label{sec:mag-phase}

\subsection{Система}

Ряд синуса обрывается на кубическом члене: $\sin\varphi\approx\varphi-\varphi^3/6$. В производной второго уравнения \eqref{eq:smooth} это замена $\cos\varphi\approx 1-\varphi^2/2$. Сначала $\Lambda=0$, и система \eqref{eq:system} принимает вид
\begin{equation}
\label{eq:system-phi}
\frac{\mathrm{d}\varphi}{\mathrm{d}t}=\omega_\rf\bigl(\eta_0 w+\tilde\eta_1 w^2\bigr),
\qquad
\frac{\mathrm{d}w}{\mathrm{d}t}=K\cos\varphi_s\Bigl(\varphi-\frac{\varphi^3}{6}\Bigr).
\end{equation}
При $\tilde\eta_1=0$ из \eqref{eq:system-phi} исключается $w$. Дифференцирование первого уравнения и подстановка второго с $\omega_\rf\eta_0 K\cos\varphi_s=-\omega_s^2$ из \eqref{eq:ws} дают
\begin{equation}
\label{eq:phi-ode}
\ddot\varphi+\omega_s^2\Bigl(\varphi-\frac{\varphi^3}{6}\Bigr)=0.
\end{equation}
Полуразмах фазы --- $\hat\varphi=\max|\varphi|$.

\subsection{Решение}

При $\tilde\eta_1=0$ решение \eqref{eq:phi-ode} получено в приложении~\ref{sec:bm-cubic}. Множитель $\operatorname{sgn}\eta_0$ уравнивает ведущий член с \eqref{eq:phi}:
\begin{align}
\varphi(t)
&=\operatorname{sgn}\eta_0\Biggl[
\hat\varphi\sin\psi
+\frac{\hat\varphi^3}{192}\bigl(\sin\psi+\sin 3\psi\bigr)
\Biggr],
\label{eq:phi-ph}\\
w(t)
&=A\Biggl[
\Bigl(1-\frac{\hat\varphi^2}{64}\Bigr)\cos\psi
+\frac{\hat\varphi^2}{64}\cos 3\psi
\Biggr],
\label{eq:w-ph}\\
\psi&=\Omega t,
\qquad
\frac{\Omega}{\omega_s}=1-\frac{\hat\varphi^2}{16}.
\label{eq:Omega-ph}
\end{align}
В ведущем порядке $\hat\varphi=\omega_\rf|\eta_0|A/\omega_s$, и \eqref{eq:phi-ph} совпадает с \eqref{eq:phi}. При $\psi=0$ и $\psi=\pi$ скобка в \eqref{eq:w-ph} равна $\pm1$, так что $A=(w_{\max}-w_{\min})/2$. Этот полуразмах связан с $\hat\varphi$ формулой
\begin{equation}
\label{eq:A-ph}
A=\frac{\omega_s\hat\varphi}{\omega_\rf|\eta_0|}\Bigl(1-\frac{\hat\varphi^2}{24}\Bigr).
\end{equation}
Для маятника $\ddot\varphi+\omega_s^2\sin\varphi=0$ интеграл энергии даёт то же соотношение точно:
\begin{equation}
\label{eq:A-exact}
A=\frac{2\omega_s}{\omega_\rf|\eta_0|}\sin\frac{\hat\varphi}{2}.
\end{equation}
Разложение \eqref{eq:A-exact} до членов $\hat\varphi^2$ совпадает с \eqref{eq:A-ph}.

При $\tilde\eta_1\neq 0$ в первом порядке по $\tilde\eta_1$ к \eqref{eq:w-ph} добавляются постоянный член и вторая гармоника \eqref{eq:w-sol}. Члены порядка $\tilde\eta_1 A^2\hat\varphi^2$ не удерживаются:
\begin{align}
w(t)
&=w_\eq
+A\Biggl[
\Bigl(1-\frac{\hat\varphi^2}{64}\Bigr)\cos\psi
+\frac{\hat\varphi^2}{64}\cos 3\psi
\Biggr]
+\frac{\tilde\eta_1 A^2}{6\eta_0}\cos 2\psi,
\label{eq:w-ph-eta}\\
\frac{\Omega}{\omega_s}
&=1-\frac{\hat\varphi^2}{16}
-\frac{5}{12}\Bigl(\frac{\tilde\eta_1 A}{\eta_0}\Bigr)^2.
\label{eq:Omega-ph-eta}
\end{align}
Постоянная $w_\eq$ определяется в \eqref{eq:deq-ph}. При $\psi=0$ и $\psi=\pi$ вторая гармоника одинакова, как в \eqref{eq:w-sol}, и $A$ остаётся полуразмахом.

\subsection{Усреднение и равновесный сдвиг}

Равенства \eqref{eq:deq-avg} точны и от второго уравнения системы не зависят: порядок по фазе входит только через $\langle w^2\rangle$. На этом порядке $\langle w^2\rangle=\langle\delta^2\rangle$. В первом порядке по $\tilde\eta_1$ среднее берётся на решении \eqref{eq:w-ph}. Гармоники $\cos\psi$ и $\cos 3\psi$ ортогональны, поэтому
\begin{equation}
\label{eq:w2-ph}
\langle w^2\rangle
=\frac{A^2}{2}\Biggl[
\Bigl(1-\frac{\hat\varphi^2}{64}\Bigr)^2
+\Bigl(\frac{\hat\varphi^2}{64}\Bigr)^2
\Biggr]
=\frac{A^2}{2}\Bigl(1-\frac{\hat\varphi^2}{32}\Bigr)+O(\hat\varphi^4).
\end{equation}
Подстановка в \eqref{eq:deq-avg} даёт
\begin{equation}
\label{eq:deq-ph}
\delta_\eq=-\frac{\eta_1 A^2}{2\eta_0}\Bigl(1-\frac{\hat\varphi^2}{32}\Bigr),
\qquad
w_\eq=-\frac{\tilde\eta_1 A^2}{2\eta_0}\Bigl(1-\frac{\hat\varphi^2}{32}\Bigr),
\qquad
\Delta E_\eq=\beta_s^2 E_s\,w_\eq .
\end{equation}

Внутри сепаратрисы для маятника $\ddot\varphi+\omega_s^2\sin\varphi=0$ при любом $\hat\varphi$
\begin{equation}
\label{eq:w2-ell}
\langle w^2\rangle
=A^2\frac{\mathrm{E}(k)/\mathrm{K}(k)-1+k^2}{k^2},
\qquad
k=\sin\frac{\hat\varphi}{2},
\end{equation}
где $\mathrm{K}$ и $\mathrm{E}$ --- полные эллиптические интегралы первого и второго рода, $k$ --- модуль. Разложение \eqref{eq:w2-ell} до членов $\hat\varphi^2$ даёт множитель $(1-\hat\varphi^2/32)$ в \eqref{eq:w2-ph}.

Поправка к $\langle w^2\rangle$ возникает потому, что $A$ есть полуразмах: кубический член рождает третью гармонику, и при фиксированном полуразмахе она перераспределяет $\langle w^2\rangle$.

\subsection{С бетатронным запаздыванием}

Первое уравнение содержит $\tau$ из \eqref{eq:tau-L}, второе --- ряд $\sin\varphi\approx\varphi-\varphi^3/6$:
\begin{equation}
\label{eq:system-phi-L}
\frac{\mathrm{d}\varphi}{\mathrm{d}t}
=\omega_\rf\bigl(\eta_0 w+\tilde\eta_1 w^2+\Lambda\bigr),
\qquad
\frac{\mathrm{d}w}{\mathrm{d}t}
=K\cos\varphi_s\Bigl(\varphi-\frac{\varphi^3}{6}\Bigr).
\end{equation}
При $\tilde\eta_1=0$ устойчивая точка лежит в $w=-\Lambda/\eta_0$. Отклонение $u=w+\Lambda/\eta_0$ удовлетворяет \eqref{eq:phi-ode}: движение около этой точки есть тот же маятник, с полуразмахом $A$ относительно смещённой точки и со связью \eqref{eq:A-ph}, \eqref{eq:A-exact}.

В первом порядке по $\tilde\eta_1$ и по $\Lambda$ величина $\langle\delta^2\rangle$ берётся из \eqref{eq:w2-ph}. Подстановка в \eqref{eq:deq-raw-L} и переход к $w$ по \eqref{eq:w-series}, \eqref{eq:tau-w} дают
\begin{equation}
\label{eq:deq-ph-L}
\begin{aligned}
\delta_\eq
&=-\frac{1}{\eta_0}\Biggl[
\frac{\eta_1 A^2}{2}\Bigl(1-\frac{\hat\varphi^2}{32}\Bigr)+\Lambda
\Biggr],
\\
w_\eq
&=\delta_\eq+\frac{A^2}{4\gamma_s^2}\Bigl(1-\frac{\hat\varphi^2}{32}\Bigr)
=-\frac{1}{\eta_0}\Biggl[
\frac{\tilde\eta_1 A^2}{2}\Bigl(1-\frac{\hat\varphi^2}{32}\Bigr)+\Lambda
\Biggr],
\\
\Delta E_\eq
&=\beta_s^2 E_s\,w_\eq .
\end{aligned}
\end{equation}
При $\Lambda=0$ это \eqref{eq:deq-ph}.

К правой части \eqref{eq:Omega-L} добавляется сдвиг $-\hat\varphi^2/16$ из \eqref{eq:Omega-ph}:
\begin{equation}
\label{eq:Omega-ph-L}
\frac{\Omega}{\omega_s}
=1-\frac{\tilde\eta_1\Lambda}{\eta_0^2}
-\frac{5}{12}\Bigl(\frac{\tilde\eta_1 A}{\eta_0}\Bigr)^2
-\frac{\hat\varphi^2}{16}.
\end{equation}

Условия применимости ряда:
\begin{equation}
\label{eq:valid-ph}
\Bigl|\frac{\eta_1 A}{\eta_0}\Bigr|\ll 1,
\qquad
\Bigl|\frac{\tilde\eta_1\Lambda}{\eta_0^2}\Bigr|\ll 1,
\qquad
\hat\varphi^2\ll 1.
\end{equation}
Формула \eqref{eq:w2-ell} верна внутри сепаратрисы при любом $\hat\varphi$.
